% DETALHE do desenho por estudo (suplemento, Secao 8). Gerado a partir do desenho original; os rotulos de secao do texto principal viram \mref{...}.
\subsection{Adaptive move selection inside HGS-CVRP (Study 7)}\label{sec:hgsaos}
Study~7 moves the question into a state-of-the-art local search. We use a modified copy of HGS-CVRP \citep{vidal2022} in which, for every pair of nodes, a policy chooses the order in which the applicable route-improvement moves (three relocations, three swaps, 2-opt and two 2-opt* variants) are tried; the first improving move is applied, as in the original, and the policy receives a binary reward (the move improved the solution). SWAP* and the insertion after the depot keep their original order. The arms are the original fixed order, a uniform random order (the non-adaptive control), UCB1, Thompson sampling, probability matching and a lite HG-UCB without the drift detector, with hyperparameters fixed in advance and not tuned (Section~7 of the online supplement). For speed, policy statistics are updated after every attempt but the indices are recomputed in batches of 16 to 256 attempts, an approximation we declare. With the original order the modified code reproduces the solutions of the original program exactly (51 of 51 runs). The budget is a number of generations (offspring), with 250 and 1000 as the primary checkpoints, read from a single run of 4000 generations because the search is deterministic given the seed. Instances are the 25 of sets X and CMT (population P1) and the 30 synthetic validation instances (P2), with ten seeds; sets A and B are excluded because HGS-CVRP converges there within a few hundred generations.

\subsection{A positive control: planted useless operators (Study 8)}\label{sec:s8design}
Null results are only informative if the protocol could have detected an effect. Study~8 therefore plants a heterogeneity of known direction. To the ten operators we add four \emph{useless} operators, each of which, when applied, consumes $X n^2$ distance evaluations and returns ``no improvement'' (it counts as a failure for the patience rule), with $X=4$ or $X=8$; for $n=100$ one application costs 40\,000 or 80\,000 evaluations, against about 13\,000 on average for a real operator (Table~\ref{tab:eff}). The levels were chosen in a pilot on the tuning instances so that uniform selection loses a measurable amount. Uniform selection and the six adaptive policies of the earlier studies run with their default hyperparameters (UCB1, probability matching and HG-UCB with the cost-aware reward, Thompson sampling, EXP3 and the efficiency UCB), on the 74 CVRPLIB instances with ten seeds at $E_1$, in the clean pool and in the two inflated pools (15\,540 runs). The confirmatory family has 12 tests (six adaptive policies $\times$ two inflated pools against uniform selection in the same pool); the share of the planted loss that a policy recovers is exploratory.

\subsection{Adaptive selection in a stochastic large-neighborhood search (Study 9)}\label{sec:s9design}
The studies above use deterministic best-improvement operators, which may be the least favorable setting for adaptation. Study~9 moves to the canonical use of adaptive selection in routing, an adaptive large neighborhood search (ALNS; \citealp{ropke2006}) whose operators are stochastic. Each iteration removes $q\sim U\{\max(5,\lceil0.1n\rceil),\min(\lceil0.3n\rceil,60)\}$ customers with one of four destroy operators (random; worst, with removal-gain rank randomization $p=3$; related, by distance and demand as in \citealp{ropke2006}, $p=6$; whole route) and reinserts them with one of three repair operators (greedy; regret-2; a cheap one that inserts customers in random order at their best position), the two choices being made by two independent instances of the same policy, as in the original method. The result is accepted by simulated annealing with initial temperature $0.02\,c_0/\ln2$ and exponential cooling driven by the fraction of the budget spent (the same for every arm), the budget being $1500\,n^2$ distance evaluations (insertion and relatedness evaluations count; the lengths of existing edges are read without counting). The configuration was fixed on the tuning instances. We use the instances with the ``at least $K$ routes'' convention, the 21 of set X and the 30 synthetic validation instances (ten and five seeds). The arms are uniform selection, the adaptive weights of \citet{ropke2006} (RP; segments of 100 iterations, scores 33, 9 and 13 for a new best, an improvement and an accepted worse solution), UCB1, probability matching, HG-UCB, Thompson sampling, EXP3 and the efficiency UCB; the hyperparameters of RP, UCB1, probability matching and HG-UCB are tuned on the 24 tuning instances by the grid protocol of Study~4 and the others keep their defaults. Two VNS arms at the same budget (uniform selection, and the static mixture of Study~4) serve as a reference for the choice of search structure. The primary test is RP against uniform selection on the 51 instances (one test); the secondary family has 12 tests (six policies, two instance sets) and the tertiary family four (ALNS with uniform selection against the two VNS arms, two sets).

\subsection{Is the gain adaptation? A static control for the ALNS (Study 12)}\label{sec:s12}
The gains of Study~9 could come from learning during the run or merely from weighting the operators better, which a fixed probability vector tuned offline also achieves. Study~12 adds two non-adaptive arms to the ALNS of Study~9: a \emph{static mixture}, with destroy probabilities $\mathrm{softmax}(\theta_{1:4})$ and repair probabilities $\mathrm{softmax}(\theta_{5:7})$ tuned on the 24 tuning instances by the $(1{+}4)$ evolution strategy of Section~\mref{sec:tuning} (eight generations, 33 evaluations of 72 runs each), and a fixed repair (always regret-2, uniform destroy). The new arms run on the same 51 instances, seeds and budget as Study~9 and are compared with the results of that study (same code). The confirmatory family has five tests (static mixture and fixed repair against uniform selection; adaptive weights against each; probability matching against the static mixture), and the central one is adaptive weights against the static mixture.

\subsection{Does the value of selection grow with instance size? (Study 10)}\label{sec:s10design}
A post hoc analysis of the data of Studies 3 and 4 suggested that the static mixture tuned on instances of up to 150 customers beats uniform selection on the larger instances (CVRPLIB instances with more than 120 customers: \num{s10ph_d_1}~pp at $E_1$, \num{s10ph_w_1} wins of \num{s10ph_n_1}, unadjusted $p=\num{s10ph_p_1}$; \num{s10ph_d_5}~pp at $E_2$, $p=\num{s10ph_p_5}$), a pattern that the growing cost of the naive SWAP* neighborhood (Table~\ref{tab:eff}) could explain. Because the subgroup was chosen after the fact, Study~10 tests it on new, larger instances. We generate two further synthetic sets with the same validated generator (and the same reference procedure): a tuning set of 12 instances (200, 300 and 400 customers, four each) and a validation set of 20 instances (200, 300, 400 and 500 customers, five each), disjoint from all other sets. A static mixture is tuned on the tuning set by the $(1{+}4)$ evolution strategy of Section~\mref{sec:tuning} (three seeds); the evaluation compares uniform selection, the mixture of Study~4 (tuned at $n\le150$), the new mixture, and the five online policies (default hyperparameters) on the validation set at $E_1$ with five seeds (800 runs). The confirmatory family has seven tests (each arm against uniform selection).

\subsection{Do the conclusions survive an efficient SWAP*? (Study 11)}\label{sec:s11design}
The SWAP* neighborhood of Studies 3--10 is evaluated naively (all pairs of customers of two routes, all insertion positions); it is much costlier in distance evaluations than the version of HGS-CVRP \citep{vidal2022} (\num{eff_cost_9}~thousand evaluations per application under uniform selection, Table~\ref{tab:eff}) and absorbs \num{s10_swap_n200}--\num{s10_swap_n500}\% of the evaluations on instances of 200 to 500 customers (Study~10). The pool effect (Study~3), the poor showing of cost-aware policies and the absence of gains from selection could therefore be artifacts of that cost. Study~11 replaces the naive SWAP* by an implementation that precomputes, for each customer and route, the three best insertion positions \citep{vidal2022}, without the polar-sector restriction: it returns the same best move as the naive version on all 36 test solutions of six instances, with 4 to 34 times (18 on average) fewer distance evaluations (unit test in the repository). Inside the search, where most applications find no improving move, the mean cost per application falls by a factor of \num{s11_cost_ratio} (Section~\ref{sec:s11}). The pool P10F is the ten-operator pool with this operator; the structure (R5 termination, best acceptance), the budgets $E_1$ and $E_2$, and the instances and seeds are those of Studies 3--6. A static mixture is tuned afresh for P10F on the 24 tuning instances by the evolution strategy of Section~\mref{sec:tuning} (\num{s11_tune_runs} runs). The evaluation compares uniform selection, six online policies (UCB1, probability matching and HG-UCB with the cost-aware reward, Thompson sampling, EXP3 and the efficiency UCB, all with default hyperparameters), the tuned mixture and, for the effect of the pool, uniform selection over the five original operators (P5), on the 74 CVRPLIB instances (ten seeds, $E_1$ and $E_2$) and, without P5, on the 20 large synthetic instances of Study~10 (five seeds, $E_1$): \num{s11_nruns} runs. The confirmatory families have 21 tests (each of the seven non-uniform arms against uniform selection in the three scenarios) and 2 tests (the pool effect at $E_1$ and $E_2$); the comparison with the naive SWAP* of Studies 3 and 6 is descriptive.

\subsection{Does it generalize beyond the CVRP? The traveling salesman problem (Study 13)}\label{sec:s13design}
Studies 1--12 use the CVRP. Study~13 repeats the apparatus on the symmetric Euclidean traveling salesman problem (TSP), which is the CVRP with one vehicle and no capacity limit. It is a modest test of generality by design: the neighborhoods are different operators written anew, but the ALNS reuses the removal and insertion machinery of Study~9. In exchange it offers an independent instance class, a near-optimal reference at low cost and none of the feasibility questions of the CVRP.

\emph{Instances and references.} Evaluation uses \num{s13_nT} EUC\_2D instances of TSPLIB \citep{reinelt1991} with \num{s13_Tmin} to \num{s13_Tmax} nodes (ten seeds; the published optimum is the reference) and \num{s13_nS} synthetic instances with \num{s13_Smin} to \num{s13_Smax} nodes, half uniform and half clustered (five seeds). A disjoint tuning set of 24 synthetic instances (50, 100 and 150 nodes) serves for tuning and calibration. The reference of a synthetic instance is the best of two independent runs of LKH-3 \citep{helsgaun2000}; the two runs agreed on 52 of 54 instances and three further runs confirmed the minimum on the other two. LKH-3 reproduces the published optimum of \num{s13_lkh_ok} of \num{s13_lkh_tot} TSPLIB instances with our cost function, which checks the reader and the rounding convention. The initial solution is the nearest-neighbor tour from node 0, the same for every arm.

\emph{Deterministic regime (VND).} The pool P9 has nine best-improvement neighborhoods: 2-opt, Or-opt(1), Swap(1,1), Or-opt(2), Or-opt(3) (these five, P5, are the analogue of the original five operators), CROSS($\le3$), Or-opt(4--10) and two 3-opt reconnections (exchange of adjacent segments, and reversal of both). A pass costs from about $N^2$ to $N^3/2$ distance evaluations, a ratio of about 60 between the cheapest and the costliest operator in the calibration. The operators are vectorized and the number of evaluations is counted analytically, as a loop implementation would count it; both were validated against an independent brute-force enumeration (11\,792 checks, and four injected faults were detected). The structure and budgets are those of Studies 4--11 (R5 termination, best acceptance, shake with the five operators of P5, $E_1=600\,n^2$ and $E_2=5E_1$ distance evaluations; an exploratory $E_3=25E_1$). At $E_1$ a run applies only about \num{s13_apps_U} operators, so the policies have little time to learn, which $E_2$ and $E_3$ mitigate. The arms are uniform selection, the six online policies of Study~11, a static mixture ($\theta\in\mathbb{R}^9$, tuned by the evolution strategy of Section~\mref{sec:tuning} on the 24 tuning instances at $E_1$: \num{s13_tune_cfg} configurations, 5 seeds) and uniform selection over P5 (the pool effect); the evaluation runs on TSPLIB at $E_1$, $E_2$ and $E_3$ and on the synthetic instances at $E_1$ (\num{s13_nruns_ev} runs). Positive control, as in Study~8: four useless operators of cost $8n^2$ (pool D8) or $16n^2$ (D16) are planted in P9 and the seven arms without the static mixture run on TSPLIB at $E_1$ (\num{s13_nruns_pc} runs).

\emph{Stochastic regime (ALNS).} Four destroy operators (random, worst, related by distance, and \emph{segment}, which removes $q$ consecutive nodes of the tour, replacing the route removal of the CVRP) and three repair operators (greedy insertion, regret-2 and cheapest insertion in random order) with the simulated-annealing acceptance and the budget of Study~9 ($1500\,n^2$ evaluations, each iteration charged $n+1$ for the evaluation of the candidate). Ten arms: uniform selection, the adaptive weights of Ropke and Pisinger, UCB1, probability matching and HG-UCB (tuned on the grids of Study~9 on the tuning instances), Thompson sampling, EXP3 and the efficiency UCB (defaults), a static mixture ($\theta\in\mathbb{R}^7$, eight generations) and a fixed repair (the best single repair operator on the tuning instances, uniform destroy), on the \num{s13a_n} evaluation instances (\num{s13_nruns_a} runs); the deterministic VND at the same budget (uniform over P9, uniform over P5 and the static mixture; \num{s13_nruns_ref} runs) is the reference for comparing the two regimes. One fact about this design was found only in the tuning runs, after the pre-registration: with a single route the regret of the second-best route is undefined, and regret-2 insertion makes the same choices as greedy insertion (the same tours in \num{s13_req_n} of \num{s13_req_tot} random partial solutions, whereas with three routes they differ in \num{s13_req3_n} of \num{s13_req3_tot}; unit test in the repository). The repair set therefore contains two distinct operators and uniform selection gives two thirds of the repairs to one of them. We kept the design, since it was pre-registered, and report the consequence in Section~\ref{sec:s13}.

\emph{Tests.} Pre-registered, with Holm correction by family: VND, 21 tests (each of the seven non-uniform arms against uniform selection in TSPLIB $E_1$, TSPLIB $E_2$ and the synthetic set) and 3 tests of the pool effect; positive control, 12 tests (six policies against uniform selection in each of the two pools); ALNS, one primary test (adaptive weights against uniform selection on the \num{s13a_n} instances), 12 secondary tests (six policies against uniform selection in each set), 4 tertiary tests (static mixture and fixed repair against uniform selection, and adaptive weights against each of them) and 6 structure tests (ALNS with uniform selection against each of the three VND references, by set). Exploratory, labeled as such: $E_3$, the contrasts of the efficiency UCB and Thompson sampling against the static mixture, and the fraction of the planted loss that each policy recovers. After seeing the primary results we added, as \emph{post hoc} analyses not in the pre-registration, a run of the classical VND with a fixed order (the order of the list above, which is that of increasing nominal cost, and its reverse; \num{s13_nruns_x} runs) and the contrasts of the best policies against it, the same three policies with the list of operators shuffled in every run (\num{s13_nruns_q} runs; a test of whether their gain depends on the initialization order), the contrasts against uniform selection over P5 and, in the ALNS, against the static mixture and the fixed repair; they are labeled post hoc wherever they appear.

\subsection{A problem with new constraints: the CVRP with time windows (Study 14)}\label{sec:s14design}
The TSP of Study~13 has neither capacity nor time constraints. Study~14 repeats the deterministic regime on the CVRP with time windows (CVRPTW), where a move that improves the cost may be infeasible in time, so that the neighborhoods must reason about arrival times, and where best-known solutions are published. It covers the descent regime only; the stochastic regime with time windows is Study~16 (Section~\ref{sec:s16design}).

\emph{Instances and cost.} Evaluation uses the \num{s14_nS} instances of \citet{solomon1987} (100 customers; families C1, C2, R1, R2, RC1 and RC2), downloaded from CVRPLIB \citep{cvrplib}, whose reference is the best known solution (proven optimal for all of them). The cost is the total distance with unlimited fleet (at most 25 vehicles), with the CVRPLIB convention that distances and travel times are truncated to one decimal, stored in tenths as integers so that cost and feasibility are exact. Our reader and feasibility simulation reproduce the published cost of the best known solution, and its feasibility, on \num{s14_bks_ok} of \num{s14_bks_tot} instances. A disjoint tuning set of 24 synthetic instances of 100 customers in the style of Solomon (four per family; windows anchored on a feasible sweep route, narrow for the type~1 families and wide for the type~2 families) serves for tuning and calibration; their reference is the best of two runs of PyVRP \citep{wouda2024} of 20~s, which agreed on 21 of 24. The initial solution is a deterministic sequential insertion by increasing due date, the same for every arm.

\emph{Operators.} The pool P9 has nine best-improvement neighborhoods: Move(1) and Swap-intra and Or-opt(2--3) within a route; Swap(1,1), Relocate(1,0), Relocate(2,0) (both orientations), 2-opt* (exchange of route tails), CROSS($\le3$) and CROSS(4--5) between routes. The first five (P5) are the analogue of the original operators. Feasibility in time is checked in constant time per candidate by concatenating precomputed segments \citep{vidal2013}. The operators are implemented in compiled code, and the number of distance evaluations is counted analytically: a candidate that passes the capacity test costs the number of new arcs of the move. Both were validated against an independent brute-force enumeration that builds every neighbor, simulates its feasibility and recomputes its cost from scratch (2\,880 checks with no failure; five injected faults were detected). Under uniform selection at $E_1$ the cost per application ranges from 2\,500 evaluations (Swap-intra) to 103\,000 (CROSS(4--5)), a factor of 42, and the two CROSS neighborhoods absorb \num{s14_cross_U}\% of the evaluations. The structure and budgets are those of Studies 4--13 ($E_1=600\,n^2$, $E_2=5E_1$, exploratory $E_3=25E_1$ distance evaluations, R5 termination, best acceptance, shake with a random feasible move of P5). At $E_1$ a run applies about \num{s14_apps_U} operators.

\emph{Arms.} Uniform selection; the six online policies and the static mixture of Study~13 (the mixture tuned by the evolution strategy of Section~\mref{sec:tuning} on the 24 tuning instances at $E_1$: \num{s14_tune_cfg} configurations, 5 seeds); uniform selection over P5; and three classical fixed-order descents without learning, pre-registered because of Study~13: the order of the list, the order of increasing cost per application measured in the calibration (Swap-intra, Move(1), 2-opt*, Or-opt(2--3), Swap(1,1), Relocate(1,0), Relocate(2,0), CROSS($\le3$), CROSS(4--5)) and its reverse. The evaluation runs on the Solomon instances at $E_1$, $E_2$ and $E_3$ with ten seeds (\num{s14_nruns_ev} runs). Positive control: four useless operators of cost $2n^2$ (pool D2) or $4n^2$ (D4), which raise the gap of uniform selection at $E_1$ in the calibration by 0.69 and 1.44~pp, are planted in P9 and the seven arms without the mixture, P5 and the orders run at $E_1$ (\num{s14_nruns_pc} runs).

\emph{Tests.} Pre-registered, with Holm correction by family: each of the seven non-uniform arms against uniform selection at $E_1$ and $E_2$ (14 tests); the pool effect (2); the three fixed orders against uniform selection at $E_1$ and $E_2$ (6); and, for the positive control, six policies against uniform selection in each pool (12). Exploratory, labeled as such: $E_3$ (11 tests), the efficiency UCB and cost-aware UCB1 and HG-UCB against the cost-ordered descent, the smaller pool and the static mixture (18 tests), and the fraction of the planted loss recovered. Any analysis not in the pre-registration is labeled post hoc where it appears.


\subsection{Is the dead weight of a pool predictive? A prospective test (Study 15)}\label{sec:s15design}
The regularity that policies gain where the pool contains operators that are not worth their cost was formulated after the data of five settings (Section~\mref{sec:r7}). Study~15 turns it into a prediction. In the TSP and in the CVRPTW the nested pools P5 $\subset$ P6 $\subset$ P7 $\subset$ P8 are prefixes of the operator list of Studies 13 and 14. The dead weight $D_k$ of the pool P$k$ is the mean gap of uniform selection over P$k$ minus that over P5, measured at $E_1$ on the 24 synthetic tuning instances of each problem (five seeds). The rule fixed before the evaluation: if $D_k\ge\num{s15_lim}$~pp at least one of the six policies (cost-aware UCB1, probability matching and HG-UCB, Thompson sampling, EXP3, the efficiency UCB, with the defaults of Studies 13 and 14) is significantly better than uniform selection (two-sided Wilcoxon on instance means, Holm over the six within the pool) on the primary set at $E_1$ (TSPLIB; Solomon); otherwise none is. The evaluation adds the fixed-order descent of each pool (list order in the TSP, order of increasing cost restricted to the pool with time windows) and, in the TSP, the 30 synthetic validation instances as an exploratory replication: \num{s15_nruns} runs. The pre-registration is log entry 118 (commit \texttt{bab11eb}).

\subsection{The stochastic regime with time windows: an ALNS for the CVRPTW (Study 16)}\label{sec:s16design}
Study~14 covers the descent regime only, which left open whether the pattern of the ALNS of Studies 9, 12 and 13 (adaptive weights beat uniform selection, but a tuned fixed design does as well) holds when insertions must respect time windows. Study~16 repeats that regime on the Solomon instances of Study~14, with the same arms, tuning protocol and budget as Study~13.

\emph{Operators.} Four destroy operators (random, worst, related and route removal) and three repair operators (greedy insertion, regret-2 and cheapest insertion in random order), chosen by two independent instances of the same policy, as in Study~9. Related removal adds a time term to the distance and demand terms of Shaw (weights 9, 3 and 2 as in \citet{ropke2006} for the VRP with time windows; the time term is the difference of the window starts). Each route stores the start of service of every node and the latest start of service that keeps the rest of the route feasible, so that the feasibility of every insertion position is decided in constant time; costs and times are integers (tenths), hence exact. The complete solution is checked again by simulation in every iteration. Because the distances are truncated to one decimal, the triangle inequality can fail by 0.1 and the removal of a customer can make a tight route infeasible; such candidates are rejected and counted (\num{s16a_inv_U} per run under uniform selection). The operators were validated against independent brute-force implementations that use only the exact simulation of a route: best insertion position and its evaluation count, the exact outcome of the three repairs, the cost bookkeeping and the feasibility of the repaired solution (9\,818 checks, no failure; nine injected faults, among them a missing latest-start test, a strict inequality, ignoring waiting and a sign error in the regret, were all detected), and the complete ALNS returned feasible solutions never below the proven optimum on five Solomon instances (unit test in the repository).

\emph{Budget and temperature.} The budget is $1500\,n^2$ distance evaluations, about 4\,100 iterations per run, each iteration charged $n+R$ evaluations for the candidate, as in Studies 9 and 13. The initial temperature $w_0$ ($T_0=w_0\,c_0/\ln 2$) mattered far more than in the TSP: in the calibration on the tuning instances (three pilots of \num{s16_pilot_runs} runs, the tuning set only) the gap of uniform selection fell from \num{s16p_U_w002}\% at $w_0=0.02$ to \num{s16p_U_w02}\% at $w_0=0.2$. The rule fixed before any data of the study, the same for every arm, was the value of the grid $\{0.02, 0.05, 0.1, 0.2, 0.3, 0.5, 1\}$ that minimizes the mean gap of uniform selection on the tuning instances, so as not to handicap the baseline with a poor temperature: $w_0=0.2$. With it, the fixed regret-2 repair was worse than uniform selection in the calibration (\num{s16p_reg2_w02}\% against \num{s16p_U_w02}\%) and the adaptive weights were slightly better (\num{s16p_RP_w02}\%), unlike in the CVRP, where regret-2 dominated the other repairs.

\emph{Arms and tuning.} Ten arms as in Study~13: uniform selection; the adaptive weights of Ropke and Pisinger, UCB1, probability matching and HG-UCB (tuned on the grids of Study~9 on the 24 tuning instances, three seeds); Thompson sampling, EXP3 and the efficiency UCB (defaults); a static mixture ($\theta\in\mathbb{R}^7$, evolution strategy of eight generations, \num{s16_tune_runs} tuning runs in all together with the grids); and a fixed repair (the best single repair on the tuning instances, uniform destroy). Evaluation: the \num{s16_n} Solomon instances with ten seeds (\num{s16_nruns_a} runs; all feasible). The deterministic VND of Study~14 at the same budget (uniform over P9 and over P5, the descent of increasing cost and the static mixture of the VND; \num{s16_nruns_ref} runs) is the reference for the comparison of the two regimes.

\emph{Tests.} Pre-registered (log entry 121, before any tuning or evaluation data of the study), with Holm correction by family: one primary test (adaptive weights against uniform selection, unadjusted); six secondary tests (the other policies against uniform selection); four tertiary tests (static mixture and fixed repair against uniform selection, adaptive weights against each); fourteen tests of each adaptive policy against each non-adaptive design (the operationalization of ``no policy beats the best fixed design''); and four structure tests (ALNS with uniform selection against each VND reference). Exploratory, labeled as such: the six Solomon classes, equivalence, and the fraction of the gain of the adaptive weights captured by the static mixture and the fixed repair.

\subsection{Adaptive selection of the operator set inside a modern solver: the ILS of PyVRP (Study 17)}\label{sec:s17design}
Study~17 answers the most frequent objection to the earlier studies, that the searches are written in Python and reach gaps of 1.3--3.5\%: it repeats the selection question inside PyVRP~0.13 \citep{wouda2024}, whose iterated local search (ILS: perturbation, granular local search and late-acceptance hill-climbing) reaches a mean gap of \num{s17p_P6_c}\% with its default operators on the tuning instances at 0.1~s per customer.

\emph{What is selected.} The local search of PyVRP scans every pair of neighboring nodes and applies all the operators of its set in turn, so there is no ``next neighborhood'' to choose as in a VND; what a policy can choose at each ILS iteration is the operator set, that is, the strength and the cost of the search. The $k=7$ arms are nested sets: P1 is Exchange(1,0), P2 adds Exchange(2,0), P3 Exchange(1,1), P4 Exchange(2,1), P5 Exchange(2,2), P6 SwapTails (this is the default set of PyVRP for the CVRP) and P7 adds the route operator SwapStar to P6. The initial solution and the exhaustive search that follows every new best solution use P6 in the adaptive arms (the arm's own set in the fixed arms). Policies, reward and tuning are those of the ALNS studies, with one decision per iteration: reward 33 for a new best feasible solution, 9 for an improvement of the current one and 0 otherwise (the acceptance rule of the ILS is not visible to the search), and the efficiency policy is charged the wall time of the call. The module is in the repository (\texttt{src/pyvrp\_aos.py}); with a single set and the same seed it reproduces \texttt{pyvrp.solve} in cost and in number of iterations on five instances and two seeds, the distances and the route costs agree with ours on every tested pair and solution, and all returned solutions satisfy our feasibility convention (unit test in the repository).

\emph{Budget.} The budget is wall-clock time, $T=\nu n$ seconds for $n$ customers, with $\nu=\num{s17_nivel_1}$ and $\num{s17_nivel_2}$ chosen, before any tuning or evaluation data, from the convergence curve of the default set on the tuning instances alone (mean gap \num{s17p_P6_a}, \num{s17p_P6_b} and \num{s17p_P6_c}\% at 0.01, 0.03 and 0.1~s per customer; the first level was discarded because a run lasts 0.5 to 1.5~s, dominated by fixed costs). Eleven processes run in parallel on twelve threads with the tasks shuffled, so the results are not reproducible bit for bit; the number of iterations and the time per set are recorded.

\emph{Instances and arms.} The 74 CVRPLIB instances (ten seeds) and the 30 synthetic validation instances (five seeds), 104 in all; the fleet is capped at $K$ for the sets A and B, which require exactly $K$ routes, and free otherwise, and runs that violate the route-count convention are excluded (\num{s17_ninv} of \num{s17_nruns}). Fourteen arms: the fixed sets P3 to P7 (P1 and P2 are almost useless and appear only in the tuning and as options of the policies), uniform selection over the seven sets, the adaptive weights of Ropke and Pisinger, UCB1, probability matching and HG-UCB (tuned on the grids of Study~9 on the 24 tuning instances, three seeds, at the first budget), Thompson sampling, EXP3 and the efficiency UCB (defaults), and a static mixture ($\theta\in\mathbb{R}^7$, evolution strategy of eight generations; \num{s17_tune_runs} tuning runs in all). The tuned parameters are used at both budgets.

\emph{Tests.} Pre-registered (log entry 124, before any tuning or evaluation data), with Holm correction by family: one primary test (adaptive weights against uniform selection at the first budget, unadjusted); fifteen secondary tests (the other policies and the static mixture against uniform selection at both budgets, and the adaptive weights at the second); eight tests of the ladder of sets (P4$-$P3 to P7$-$P6 at both budgets); four tertiary tests (static mixture against P6 and P7); two tests of uniform selection against the default set; and 21 tests per budget of each adaptive policy against P6, P7 and the static mixture (the operationalization of ``no policy beats the best non-adaptive design''). Exploratory, labeled as such: the two instance sets and the four families of CVRPLIB, and the usage of the sets.
